% Options for packages loaded elsewhere
% Options for packages loaded elsewhere
\PassOptionsToPackage{unicode}{hyperref}
\PassOptionsToPackage{hyphens}{url}
\PassOptionsToPackage{dvipsnames,svgnames,x11names}{xcolor}
%
\documentclass[
]{report}
\usepackage{xcolor}
\usepackage[top=30mm,left=20mm,right=20mm,bottom=30mm]{geometry}
\usepackage{amsmath,amssymb}
\setcounter{secnumdepth}{5}
\usepackage{iftex}
\ifPDFTeX
  \usepackage[T1]{fontenc}
  \usepackage[utf8]{inputenc}
  \usepackage{textcomp} % provide euro and other symbols
\else % if luatex or xetex
  \usepackage{unicode-math} % this also loads fontspec
  \defaultfontfeatures{Scale=MatchLowercase}
  \defaultfontfeatures[\rmfamily]{Ligatures=TeX,Scale=1}
\fi
\usepackage{lmodern}
\ifPDFTeX\else
  % xetex/luatex font selection
\fi
% Use upquote if available, for straight quotes in verbatim environments
\IfFileExists{upquote.sty}{\usepackage{upquote}}{}
\IfFileExists{microtype.sty}{% use microtype if available
  \usepackage[]{microtype}
  \UseMicrotypeSet[protrusion]{basicmath} % disable protrusion for tt fonts
}{}
\makeatletter
\@ifundefined{KOMAClassName}{% if non-KOMA class
  \IfFileExists{parskip.sty}{%
    \usepackage{parskip}
  }{% else
    \setlength{\parindent}{0pt}
    \setlength{\parskip}{6pt plus 2pt minus 1pt}}
}{% if KOMA class
  \KOMAoptions{parskip=half}}
\makeatother
% Make \paragraph and \subparagraph free-standing
\makeatletter
\ifx\paragraph\undefined\else
  \let\oldparagraph\paragraph
  \renewcommand{\paragraph}{
    \@ifstar
      \xxxParagraphStar
      \xxxParagraphNoStar
  }
  \newcommand{\xxxParagraphStar}[1]{\oldparagraph*{#1}\mbox{}}
  \newcommand{\xxxParagraphNoStar}[1]{\oldparagraph{#1}\mbox{}}
\fi
\ifx\subparagraph\undefined\else
  \let\oldsubparagraph\subparagraph
  \renewcommand{\subparagraph}{
    \@ifstar
      \xxxSubParagraphStar
      \xxxSubParagraphNoStar
  }
  \newcommand{\xxxSubParagraphStar}[1]{\oldsubparagraph*{#1}\mbox{}}
  \newcommand{\xxxSubParagraphNoStar}[1]{\oldsubparagraph{#1}\mbox{}}
\fi
\makeatother


\usepackage{longtable,booktabs,array}
\usepackage{calc} % for calculating minipage widths
% Correct order of tables after \paragraph or \subparagraph
\usepackage{etoolbox}
\makeatletter
\patchcmd\longtable{\par}{\if@noskipsec\mbox{}\fi\par}{}{}
\makeatother
% Allow footnotes in longtable head/foot
\IfFileExists{footnotehyper.sty}{\usepackage{footnotehyper}}{\usepackage{footnote}}
\makesavenoteenv{longtable}
\usepackage{graphicx}
\makeatletter
\newsavebox\pandoc@box
\newcommand*\pandocbounded[1]{% scales image to fit in text height/width
  \sbox\pandoc@box{#1}%
  \Gscale@div\@tempa{\textheight}{\dimexpr\ht\pandoc@box+\dp\pandoc@box\relax}%
  \Gscale@div\@tempb{\linewidth}{\wd\pandoc@box}%
  \ifdim\@tempb\p@<\@tempa\p@\let\@tempa\@tempb\fi% select the smaller of both
  \ifdim\@tempa\p@<\p@\scalebox{\@tempa}{\usebox\pandoc@box}%
  \else\usebox{\pandoc@box}%
  \fi%
}
% Set default figure placement to htbp
\def\fps@figure{htbp}
\makeatother





\setlength{\emergencystretch}{3em} % prevent overfull lines

\providecommand{\tightlist}{%
  \setlength{\itemsep}{0pt}\setlength{\parskip}{0pt}}



 


\makeatletter
\@ifpackageloaded{caption}{}{\usepackage{caption}}
\AtBeginDocument{%
\ifdefined\contentsname
  \renewcommand*\contentsname{Table of contents}
\else
  \newcommand\contentsname{Table of contents}
\fi
\ifdefined\listfigurename
  \renewcommand*\listfigurename{List of Figures}
\else
  \newcommand\listfigurename{List of Figures}
\fi
\ifdefined\listtablename
  \renewcommand*\listtablename{List of Tables}
\else
  \newcommand\listtablename{List of Tables}
\fi
\ifdefined\figurename
  \renewcommand*\figurename{Figure}
\else
  \newcommand\figurename{Figure}
\fi
\ifdefined\tablename
  \renewcommand*\tablename{Table}
\else
  \newcommand\tablename{Table}
\fi
}
\@ifpackageloaded{float}{}{\usepackage{float}}
\floatstyle{ruled}
\@ifundefined{c@chapter}{\newfloat{codelisting}{h}{lop}}{\newfloat{codelisting}{h}{lop}[chapter]}
\floatname{codelisting}{Listing}
\newcommand*\listoflistings{\listof{codelisting}{List of Listings}}
\makeatother
\makeatletter
\makeatother
\makeatletter
\@ifpackageloaded{caption}{}{\usepackage{caption}}
\@ifpackageloaded{subcaption}{}{\usepackage{subcaption}}
\makeatother
\usepackage{bookmark}
\IfFileExists{xurl.sty}{\usepackage{xurl}}{} % add URL line breaks if available
\urlstyle{same}
\hypersetup{
  colorlinks=true,
  linkcolor={blue},
  filecolor={Maroon},
  citecolor={Blue},
  urlcolor={Blue},
  pdfcreator={LaTeX via pandoc}}


\author{}
\date{}
\begin{document}

\renewcommand*\contentsname{Table of contents}
{
\hypersetup{linkcolor=}
\setcounter{tocdepth}{2}
\tableofcontents
}

\chapter{PROJECT\_SPEC:
cloudpedagogy-research-object-template}\label{project_spec-cloudpedagogy-research-object-template}

\section{1. Repo Name}\label{repo-name}

\texttt{cloudpedagogy-research-object-template}

\section{2. One-Sentence Purpose}\label{one-sentence-purpose}

A standardized template for authoring research outputs and technical
documentation with embedded AI governance metadata.

\section{3. Problem the App Solves}\label{problem-the-app-solves}

The lack of standardized transparency in AI-assisted research outputs;
provides a consistent structure for tracking and declaring AI usage and
governance completion.

\section{4. Primary User / Audience}\label{primary-user-audience}

Researchers, postgraduate students, and technical writers within the
CloudPedagogy ecosystem.

\section{5. Core Role in the CloudPedagogy
Ecosystem}\label{core-role-in-the-cloudpedagogy-ecosystem}

The ``Publication Layer''; used to generate the human-readable research
findings, technical guides, and reports produced by the other ecosystem
tools.

\section{6. Main Entities / Data
Structures}\label{main-entities-data-structures}

\begin{itemize}
\tightlist
\item
  \textbf{ProjectMetadata}: YAML frontmatter (tracking Title, Version,
  \texttt{ai\_used}, \texttt{governance\_complete}).
\item
  \textbf{Research Object}: The final bundled output (HTML/PDF)
  representing a ``Research Object'' in the ecosystem.
\end{itemize}

\section{7. Main User Workflows}\label{main-user-workflows}

\begin{enumerate}
\def\labelenumi{\arabic{enumi}.}
\tightlist
\item
  \textbf{Initialize Project}: Scaffold a new documentation or research
  repository from this template.
\item
  \textbf{Configure Metadata}: Complete the
  \texttt{project-metadata.yml} declaration, specifically tracking AI
  involvement.
\item
  \textbf{Author Content}: Write in Markdown (.qmd) using the
  pre-defined SCSS/styling system.
\item
  \textbf{Publish}: Build and deploy the static result as a governed
  research object.
\end{enumerate}

\section{8. Current Features}\label{current-features}

\begin{itemize}
\tightlist
\item
  Integrated Quarto build support.
\item
  Specialized SCSS styling (CloudPedagogy theme).
\item
  Lightweight ``Governance Ledger'' via YAML-based declaration.
\item
  Support for static HTML and PDF output formats.
\end{itemize}

\section{9. Stubbed / Partial / Incomplete
Features}\label{stubbed-partial-incomplete-features}

\begin{itemize}
\tightlist
\item
  Not explicitly defined.
\end{itemize}

\section{10. Import / Export and Storage
Model}\label{import-export-and-storage-model}

\begin{itemize}
\tightlist
\item
  \textbf{Storage}: File-system based (.yml, .qmd).
\item
  \textbf{Output}: Static files generated via the build process.
\end{itemize}

\section{11. Relationship to Other CloudPedagogy
Apps}\label{relationship-to-other-cloudpedagogy-apps}

Acts as the final ``Display'' or ``Communication'' layer for findings
generated by the mapping, simulation, and governance tools.

\section{12. Potential Overlap or Duplication
Risks}\label{potential-overlap-or-duplication-risks}

Overlaps with generic academic templates; distinguished by its metadata
fields geared specifically towards AI-governance declarations.

\section{13. Distinctive Value of This
App}\label{distinctive-value-of-this-app}

Treats ``AI Declaration'' as a first-class citizen of the research
object's metadata, making it easier to track AI-involvement across a
library of publications.

\section{14. Recommended Future
Enhancements}\label{recommended-future-enhancements}

(Inferred) Automated population of the metadata YAML from
\texttt{decision-record} or \texttt{governance-scanner} JSON exports;
bibtex-aware AI attribution.

\section{15. Anything Unclear or Inferred from Repo
Contents}\label{anything-unclear-or-inferred-from-repo-contents}

Use of Quarto as the primary engine is central to the project's build
logic and file structure.




\end{document}
